\section{Might, Cunning, Justice}

In the October 2014 issue of Culture Wars\footnote{\url{https://www.culturewars.com/}}, there is an extensive review by Thomas Paleologos of the HBO series The Wire\footnote{\url{https://en.wikipedia.org/wiki/The\_Wire}}. Homer's two epic poems illustrate two dissimilar social organizing principles, both opposed to Justice. The \emph{Iliad} is the epic of might (\emph{ble}), the notion that ``might makes right''. The \emph{Odyssey} is the epic of cunning intelligence (\emph{metis}), i.e., the notion that intelligence is merely instrumental and its purpose is entirely immanent to the world. The transcendent organizing principle, represented by Zeus, is justice (\emph{dike}).

In \emph{The Wire}, on the other hand, there is no justice, only might and cunning. The world of the streets and drug dealers is ruled by might or \emph{ble}. Force and violence are used indiscriminately to achieve desired ends.

Opposed to that world is the world of bureaucracy and politics, ruled by cunning or \emph{metis}. The nominal purpose of government and the police department is to enforce justice against the world of might. However, in order to do that, a system of bureaucracy is created, the effect of which is the creation of an alternate reality, an artificial creation of the human mind. The participants, therefore, must employ cunning in order to manipulate, and even survive, within that system. Hence, it is based on lies, deceit, and manipulations.

\textbf{Homer}, the ``teacher of all Greeks'', is also our teacher, we who are the heirs to that civilization. Justice, even when ignored, cannot be evaded and will appear as the natural law and fate. An action will set in motion a series of events, the results of which are seldom what the agents desired or intended. Not even God Almighty will interfere with that process. Therefore, Mr. Paleologos concludes that the paths of might and cunning lead to death. In an allegedly democratic regime, the populace is responsible for its own fate through its voting habits.

\paragraph{Ferguson}


\begin{quotex}
nothing is harder or more unfair than human reality: humans live in a world where it's words and not deeds that have power, where the ultimate skill is mastery of language. This is a terrible thing because basically we are primates who've been programmed to eat, sleep, reproduce, conquer and make our territory safe, and the ones who are most gifted at that, the most animal types among us, always get screwed by the others, the fine talkers, despite these latter being incapable of defending their own garden or bringing a rabbit home for dinner or procreating properly. Humans live in a world where the weak are dominant.
\flright{\textsc{Muriel Barbery}, \emph{The Elegance of the Hedgehog}}
\end{quotex}


Coincidently, recent events in Ferguson, Missouri illustrate this theme. The streets are apparently ruled by might, since youths can enter commercial establishments and simply take things without any fear of retribution. Last month, one such youth was gunned down in the street by a police officer, triggering several nights of riots, which are likewise part of the world of might. The question of justice, i.e., whether the youth's death was just or not, is still to be officially decided.

Now in the second episode of \emph{The Wire} there is an instance of unprovoked police brutality. (The viewer is the all-seeing eye who knows the truth.) The Captain meets with the cops and tells them how to ``spin'' and coordinate their stories before the internal investigation begins. Thus, words are what count, not the actual deeds. For the moiety that receives its worldview from popular culture such as this, then no appeal to ``facts'' will overcome the impression made. The other moiety persists on being surprised by this.

\paragraph{Picking up Apples off the Ground}


Recently, a fellow wrote to me that he respected my ``work'' but disagreed with some of my conclusions. Actually he probably disagrees with the premises, for otherwise he would have to accept the conclusions. There is the common practice of those who reject the modern world to cook up a Mulligan stew of various thinkers, seldom compatible with each other. The end result is not tasty intellectual fare, but rather mental indigestion.

They may collect various ``conclusions'' of someone like \textbf{Rene Guenon} much like a passerby who gathers apples off the ground. No matter how delicious is the apple pie made from them, the real value lies in wrestling with the premises of first rate thinkers, not in accepting their conclusions for purposes for which they were never intended.

Rare is the genius like \textbf{Isaac Newton} who will look at the apples on the ground and wonder how they got there.

\paragraph{Critical Theory}


\begin{quotex}
Philosophers have only interpreted the world in certain ways; the point is to change it.
\flright{\textsc{Karl Marx}, \emph{11th of the Theses on Feuerbach}}
\end{quotex}


Theories abound about the creation of the modern world: liberalism, Christianity, etc. Bracketing out, for now, the fact that such vague terms do not really refer to unitary ideas, the point is that ideas in themselves have no causative power. That is because an idea is a potentiality which cannot bring anything into existence. That would require an agent, i.e., a person who holds an idea in his intellect. Since the will follows the intellect, then we can say the idea has some power.

Most often, something called cultural Marxism is regarded as the cause. This is more properly called ``critical theory'' since it encompasses a much wider range of ideas. Although it takes some effort to master, it does leave a lot of attractive apples on the ground for others to gather. This is the educated moiety: they don't get their worldview from popular culture. They are intelligent enough to grasp the main conclusions of critical theory, yet are not able to understand how those conclusions were derived. The Stanford Encyclopedia of Philosophy\footnote{\url{https://plato.stanford.edu/entries/critical-theory/}} describes the fundamental criteria of a critical theory:

\begin{quotex}
A critical theory is adequate only if it meets three criteria: it must be explanatory, practical, and normative, all at the same time. That is, it must

\begin{enumerate}


\item explain what is wrong with current social reality

\item identify the actors to change it

\item provide both clear norms for criticism and achievable practical goals for social transformation
\end{enumerate}

Any truly critical theory of society ... has as its object human beings as producers of their own historical form of life.
\end{quotex}


It is pointless to enter into debate with a critical theory over its alleged contradictions or hypocrisies. Rather, an alternative critical theory is called for. Rather than the foundation of Marxism and psychoanalysis, it would be based on Traditional metaphysical principles. Only then can the three criteria be met. Despite all the Internet writings and international conferences, there is no agreement on these criteria. What follows are some possibilities:

\begin{enumerate}


\item Explain how current social reality is not compatible with Traditional notions of justice and an organic relationship among the various functional groups in society. Explain how the loss of Being results in nihilism, despair, drug abuse, etc.

\item The actors have been identified by Rene Guenon in The Crisis of the Modern World. If the modern world is based on democracy, then Tradition cannot likewise be so.

\item Social transformation needs to start with the organizations themselves.
\end{enumerate}


Critical theory focuses on the processes by which political consciousness is created. For example, these three are mention in Wikipedia\footnote{\url{http://en.wikipedia.org/wiki/Critical\_theory}}:

\begin{itemize}


\item The semiotic rules by which objects obtain symbolic meanings (Barthes).

\item The psychological processes by which the phenomena of everyday consciousness are generated (psychoanalytic thinkers).

\item The episteme that underlies our cognitive formations (Foucault).
\end{itemize}


These are secular apes of Tradition. Specifically,

\begin{itemize}


\item Tradition identifies the real meaning behind religious and political symbolism

\item Esoteric psychology has a deeper understanding of the human psyche than its psychological counterfeits.

\item Over the last few months, and in different contexts, we have identified many errors of conception that maintain people in illusion and false consciousness. We have also pointed out the Traditional teachings on achieving knowledge.
\end{itemize}



\flrightit{Posted on 2014-10-27 by Cologero}

\begin{center}* * *\end{center}

\begin{footnotesize}
\begin{sffamily}

\texttt{Jacob on 2014-10-29 at 16:42 said:}

I hate to sound condescending, but isn't that exactly what Cologero has written against? Being caught up in meaningless assertions of the world ``he's guilty... ... .he's not guilty''. Justice will ultimately flow from above regardless of human opinion. Anyone feel free to correct me if I'm wrong.

With that said Cologero, I think this is a great post. I do have to ask though; do you believe that using the tools of critical theory to restore the discipline of philosophy is a futile task or not.

\hfill

\texttt{Cologero on 2014-10-30 at 10:12 said:}

That's a good segue, Jacob, for upcoming posts. Critical theory is the opposite of Tradition, yet it comprises a comprehensive philosophical system with conclusions that are appealing to many today. Hence, it needs to be understood radically, at its roots, to wrestle with those who create critical thought, and not the multitude of passive consumers of ideas. Joseph de Maistre famously wrote that a revolution in the opposite direction is insufficient. The 20th century is proof enough. Also, it is pointless to debate the revolution on the foundations of science or rationality, since the understanding of rationality is fundamentally different. Instead, he claimed, what is required is the opposite of a revolution. I don't believe there has ever been the opposite of a revolution, so there needs to be some speculation about what it will look like.

\hfill

\texttt{Alistair Fraser on 2016-10-29 at 20:38 said:}

The film called ``spartacus'' is a good example of the expanding dynamics of a ``Revolution'' where there was considered to be a just cause , I'm not completely sure of the motivations of the so called character spartacus beyond some type of anti-slavery movement , but the kubrick film served out a neat observation , for most of the film, the spartacus character is the great leader of a revolution to great acclaim etc , but further on in the film when the spartacus revolution is brought to a halt , there is the scene where they are trying to establish the identity of the spartacus character in which one of his men makes the infamous statement ``I'm spartacus'' which is followed by all the men shouting out I'm spartacus , and this moment can be shown as the distinct moment in time where their physical defeat in revolution then transcended into a metaphysical quality of self-realisation through revelation , i.e. the ``I'm spartacus'' idea arrived via a higher impulse which led to revelatory connotations in transcendence of mens character , as now the revolution had transcended its need for a pinnacle point physical leader and each of the men were a symbol of the so called revolution which the spartacus was initially the sole representative . probably explains why hitler still gets massive negative press today whilst stalin and other participants in huge physical wars get little , the adolf must still be perceived as an ongoing potential metaphysical revelatory danger, by the modern medes who control revelations

\hfill

\texttt{Sigurd on 2017-08-25 at 07:23 said:}
Yes, I'm of a similar position on the matter: critical theory needs to be examined for what it is from the lens of tradition. I'm in an interesting spot because I find the subject of semiotics and social ritual a la Baudrillard's Simulacra and the notorious ``Society of the Spectacle'' to have a more-or-less adequate examination of some social phenomenon, but I nevertheless can recognize the ``apeness'' of them, due to a passive attention to the interesting chaos magician/post-Marxian crowd. I became aware of this once I read in Algis Uzdavinys work a relatively small chapter on traditional semiotics in the form of symbols and rituals. His citing of Baudrillard as the adequate conclusion of a society without tradition was rather fascinating.

As is known, semiology is quite intimate with the field of modern linguistics, particularly Chomskyan, that more or less proceeds from a phenomenology (Husserl), notwithstanding their attempt to slough off idealism in favor of post-Marxian dialectico-structural plumage. Do you know of any works examining critical theory and perhaps Husserlian phenomenology from the perspective of Traditional semiology? Really, any works elaborating on traditional symbols and ritual would be of great importance to me.

Further, is there a forum or a way I can get in contact with the members of this site? I would love to have more immediate discussion.


\hfill

\end{sffamily}
\end{footnotesize}

